\subsection{Методы второго порядка}
\label{subsec:second-order}

Будем считать, что $f(x)$ строго выпукла и дважды непрерывно дифференцируема,
так что существует гессиан $H_f(x)$. Разложим $f$ в окрестности точки $y$:
\[
  f(x) = \underbrace{f(y) + \nabla\T f(y)\,(x - y)
         + \tfrac{1}{2}(x-y)\T H(y)\,(x-y)}_{\psi(x)} + \dots
\]
Отбрасывая члены выше второго порядка, получаем
\[
  f(x) - f(y) \ \approx\ \psi(x) - f(y) ,
\]
где $\psi(x)$ — \textbf{квадратичная аппроксимация} функции $f$.

Минимизируем приближение по приращению $x - y$:
\[
  \frac{\partial\big(\psi(x) - f(y)\big)}{\partial (x-y)}
  = \nabla\T f(y) + H(y)\,(x - y) = 0
  \qquad\Longrightarrow\qquad
  x - y = -H^{-1}(y)\,\nabla f(y) ,
\]
откуда итерационная схема
\begin{equation}\label{eq:newton}
  x_{k+1} = x_k - H^{-1}(x_k)\, \nabla f(x_k) .
\end{equation}

\subsubsection*{Идея методов второго порядка}

\[
  x_0 \ \longrightarrow\ x_k :
  \quad
  \begin{cases}
    H(x_k) \ \longrightarrow\ H^{-1}(x_k), \\
    \nabla f(x_k)
  \end{cases}
  \quad \Longrightarrow \quad
  x_{k+1} = x_k - H^{-1}(x_k)\,\nabla f(x_k) .
\]

Если оставить схему~\eqref{eq:newton} как есть — это \textbf{метод Ньютона}.
При усилении
\[
  x_{k+1} = x_k - \alpha_k H^{-1}(x_k)\, \nabla f(x_k)
\]
получаем различные методы (выбор $\alpha_k$ — как в градиентных методах:
постоянный шаг, дробление, наискорейший спуск).

Шаг можно выбирать в виде
\[
  p_k = -\alpha_k H^{-1}(x_k)\,\nabla f(x_k)
  \qquad\Longleftrightarrow\qquad
  H(x_k)\, p_k = -\alpha_k \nabla f(x_k) ,
\]
то есть на каждой итерации либо обращение матрицы, либо решение СЛАУ.

\begin{idea}
Форма $H(x_k)p_k = -\alpha_k \nabla f(x_k)$ практически всегда предпочтительнее:
обращение матрицы стоит порядка $n^3$ операций и численно неустойчиво, а решение
системы с симметричной положительно определённой матрицей делается разложением
Холецкого — тот же порядок, но меньше констант и лучше точность. При больших $n$
систему решают приближённо, итерационно, не формируя $H$ целиком.

Связь с \S\ref{subsec:descent-direction}: здесь $F(x_k) = H^{-1}(x_k)$, и
положительная определённость $F$ (а значит корректность направления спуска)
следует из строгой выпуклости $f$. Если $f$ не выпукла, гессиан может оказаться
неположительно определённым — тогда ньютоновское направление перестаёт быть
направлением спуска, и матрицу приходится регуляризовать.

Выигрыш методов второго порядка — в том, что оценки~\eqref{eq:q-min} и
теоремы~\ref{thm:const-step}, упирающиеся в отношение $m/M$, здесь не возникают:
квадратичную функцию схема~\eqref{eq:newton} решает за один шаг при любой
обусловленности, а вблизи минимума даёт квадратичную скорость сходимости
($p = 2$ в смысле определения из~\S\ref{subsec:multidim-basics}). Плата —
вычисление и хранение $H$ на каждой итерации.
\end{idea}
